\chapter{Nephrectomy}

\section*{Introduction \& Clinical Indications}
A nephrectomy is a major surgical procedure involving the removal of a kidney, or a portion thereof. This can be performed as a radical nephrectomy, which entails the excision of the entire kidney, often along with the adrenal gland, perirenal fat, and Gerota's fascia, typically for larger or more aggressive renal cancers. Alternatively, a partial nephrectomy, also known as nephron-sparing surgery, involves removing only the diseased portion of the kidney, such as a tumor, while preserving as much healthy renal parenchyma as possible. The procedure is primarily indicated for renal malignancies, severe non-functioning kidneys due to benign conditions, or for living kidney donation. It can be performed through open surgery, laparoscopically, or with robotic assistance, each offering distinct advantages depending on the clinical scenario and surgeon's expertise.

\begin{itemize}
  \item Kidney removal surgery
  \item Renal excision
  \item Nephron-sparing surgery (NSS)
  \item Laparoscopic kidney removal
  \item Robotic kidney surgery
\end{itemize}

The fundamental principle of nephrectomy is the complete and safe removal of a diseased or non-functional kidney, or a portion thereof, while minimizing morbidity and preserving overall renal function when feasible. For radical nephrectomy, the primary goal is oncologic control, achieved by excising the entire kidney and its surrounding structures within Gerota's fascia to ensure clear surgical margins and prevent local recurrence. This involves meticulous dissection to isolate and ligate the renal artery and vein, followed by division of the ureter and removal of the specimen. The rationale is to eradicate the primary tumor and any potential local spread, offering the best chance for cure in localized renal malignancies.

In contrast, partial nephrectomy aims to achieve equivalent oncologic outcomes for localized renal masses while maximizing the preservation of healthy renal tissue. This is crucial for patients with pre-existing chronic kidney disease, a solitary kidney, or bilateral tumors, as it mitigates the risk of developing new-onset or worsening chronic kidney disease. The technical principle involves:
\begin{itemize}
  \item Precise localization of the tumor, often aided by intraoperative ultrasound.
  \item Temporary clamping of the renal artery (and sometimes vein) to create a bloodless field, minimizing warm ischemia time to protect the remaining kidney parenchyma.
  \item Sharp excision of the tumor with a small rim of healthy tissue (negative margins).
  \item Meticulous reconstruction of the renal collecting system and parenchyma (renorrhaphy) using sutures to prevent urinary leakage and control bleeding.
\end{itemize}
Both approaches require careful attention to vascular control and avoidance of injury to adjacent organs, guided by a thorough understanding of renal and perirenal anatomy.

The history of nephrectomy is marked by significant advancements in surgical technique, anesthesia, and understanding of renal physiology. The first successful nephrectomy was performed by Gustav Simon in Heidelberg, Germany, in 1869. His pioneering work on a patient suffering from a ureterovaginal fistula demonstrated that humans could survive and thrive with a single kidney, paving the way for future renal surgery. Early open nephrectomies were often fraught with high mortality rates due to infection, hemorrhage, and lack of adequate anesthesia.

Throughout the early 20th century, improvements in aseptic technique, blood transfusion, and anesthetic agents made nephrectomy a safer and more routine procedure. The development of specialized surgical instruments and a deeper understanding of renal anatomy further refined open surgical approaches. A major paradigm shift occurred in 1990 when Ralph Clayman and colleagues performed the first laparoscopic nephrectomy, ushering in the era of minimally invasive renal surgery. This technique offered reduced pain, shorter hospital stays, and quicker recovery times compared to traditional open surgery. The early 21st century witnessed the widespread adoption of robotic-assisted laparoscopic nephrectomy, which further enhanced surgical precision, dexterity, and visualization, particularly for complex partial nephrectomies, solidifying its role in modern urologic practice.

\begin{itemize}
  \item To treat localized renal cell carcinoma (RCC) or other kidney cancers, where surgical removal offers the best chance for cure.
  \item For cytoreductive purposes in metastatic RCC, to reduce tumor burden and improve the efficacy of systemic therapies.
  \item To remove a severely damaged or non-functioning kidney due to chronic infection (e.g., pyelonephritis), severe hydronephrosis, or extensive renal stone disease causing intractable pain or recurrent sepsis.
  \item To manage severe renal trauma that results in an irreparable kidney injury or uncontrolled hemorrhage.
  \item As a procedure for living kidney donation, providing a healthy kidney for transplantation to a recipient with end-stage renal disease.
  \item To alleviate symptoms from benign but large or symptomatic renal masses, such as angiomyolipomas or large cysts, that are causing pain or bleeding.
  \item In rare cases, for renovascular hypertension caused by a severely diseased kidney that is not amenable to revascularization.
  \item To create space or alleviate symptoms in patients with massively enlarged polycystic kidneys prior to or after renal transplantation.
\end{itemize}

Nephrectomy serves as both a primary and, in certain contexts, a secondary or adjunctive procedure. For localized renal cell carcinoma, it is typically considered the primary curative treatment, especially for T1 and T2 tumors. Partial nephrectomy is often the primary choice for smaller renal masses (T1a and T1b) due to its ability to preserve renal function without compromising oncologic outcomes. Similarly, for a severely diseased, non-functioning, or symptomatic benign kidney, nephrectomy is the primary intervention to resolve the underlying problem.

However, nephrectomy can also be a secondary or adjunctive procedure. In cases of metastatic renal cell carcinoma, a cytoreductive nephrectomy may be performed after or in conjunction with systemic therapies (e.g., immunotherapy, targeted therapy) to improve overall survival, manage symptoms like pain or bleeding, or reduce tumor burden. It may also be considered a secondary procedure if initial non-surgical treatments have failed or if a patient develops complications requiring surgical intervention. For living kidney donation, it is a primary elective procedure performed on a healthy individual for the benefit of another.

\section*{Relevant Surgical Anatomy}
The kidneys are paired retroperitoneal organs situated on either side of the vertebral column, typically extending from the level of the T12 to L3 vertebrae. The right kidney is usually slightly lower than the left due to the presence of the liver. Each kidney is encased within Gerota's fascia (also known as the renal fascia), a fibrous capsule that encloses the kidney, adrenal gland, and perirenal fat, providing a natural plane for surgical dissection, especially during radical nephrectomy.

Key anatomical relations include:
\begin{itemize}
  \item \textbf{Superiorly:} The adrenal (suprarenal) gland sits atop each kidney, often removed during radical nephrectomy for malignancy. The diaphragm is also superior.
  \item \textbf{Anteriorly:} The right kidney is anteriorly related to the liver, duodenum, and hepatic flexure of the colon. The left kidney is anteriorly related to the spleen, pancreas, stomach, jejunum, and splenic flexure of the colon.
  \item \textbf{Posteriorly:} Both kidneys lie on the psoas major, quadratus lumborum, and transversus abdominis muscles. The subcostal, iliohypogastric, and ilioinguinal nerves run posterior to the kidney and are susceptible to injury during flank incisions.
  \item \textbf{Vascular Supply:} The renal arteries, typically arising directly from the aorta at the level of L1-L2, are usually single but can have accessory branches. They branch into segmental arteries within the renal hilum. The renal veins drain into the inferior vena cava (IVC); the left renal vein is longer, passing anterior to the aorta and posterior to the superior mesenteric artery, and receives the left gonadal and left adrenal veins. The right renal vein is shorter and drains directly into the IVC.
  \item \textbf{Ureter:} The ureter originates from the renal pelvis at the hilum, descends retroperitoneally, passing anterior to the psoas muscle and posterior to the gonadal vessels, before entering the bladder.
  \item \textbf{Lymphatic Drainage:} Lymphatics from the kidneys drain primarily to the para-aortic and paracaval lymph nodes, located around the great vessels.
\end{itemize}
Understanding these intricate relationships is paramount for safe and effective nephrectomy, particularly in minimizing injury to adjacent organs and ensuring complete tumor removal.

\section*{Preoperative Workup \& Preparation}
Thorough preoperative preparation is essential to ensure patient safety and optimize surgical outcomes. This typically involves a comprehensive medical evaluation to assess overall health and identify any potential risks.

Key preparatory steps include:
\begin{itemize}
  \item \textbf{Medical History and Physical Exam:} A detailed review of the patient's medical history, including comorbidities, medications, allergies, and prior surgeries.
  \item \textbf{Laboratory Tests:}
    \begin{itemize}
      \item Complete blood count (CBC) to check for anemia or infection.
      \item Comprehensive metabolic panel (CMP) to assess electrolyte balance, liver function, and crucial renal function (serum creatinine, estimated glomerular filtration rate - GFR) of both kidneys.
      \item Coagulation profile (PT/INR, PTT) to evaluate bleeding risk.
      \item Urinalysis and urine culture to rule out urinary tract infection.
      \item Blood type and crossmatch for potential blood transfusion.
    \end{itemize}
  \item \textbf{Imaging Studies:}
    \begin{itemize}
      \item Computed tomography (CT) scan of the abdomen and pelvis with intravenous contrast is the primary imaging modality for diagnosis, staging of renal masses, and surgical planning.
      \item Chest X-ray or CT chest to assess lung health and rule out metastatic disease.
      \item Bone scan or MRI may be ordered if there is suspicion of bone metastases.
    \end{itemize}
  \item \textbf{Cardiac and Pulmonary Clearance:} Evaluation by a cardiologist or pulmonologist may be necessary for patients with significant cardiac or respiratory comorbidities to ensure they are fit for surgery and anesthesia.
  \item \textbf{Medication Management:} Instructions on which medications to stop (e.g., blood thinners, certain diabetes medications) or adjust before surgery.
  \item \textbf{NPO Status:} Patients are instructed to be NPO (nil per os - nothing by mouth) for a specified period (typically 6-8 hours) before surgery to prevent aspiration during anesthesia.
  \item \textbf{Bowel Preparation:} In some cases, particularly for open or complex procedures, a bowel preparation may be prescribed.
  \item \textbf{Pre-operative Antibiotics:} Administered intravenously shortly before the incision to reduce the risk of surgical site infection.
  \item \textbf{Informed Consent:} Detailed discussion with the surgeon about the procedure, potential risks, benefits, and alternatives, followed by signing of the informed consent form.
\end{itemize}

Contraindications to nephrectomy can be absolute or relative, depending on the patient's overall health, the nature of the renal disease, and the availability of alternative treatments.

\textbf{Situations requiring renal-risk assessment, optimization, or an alternative treatment plan:}
\begin{itemize}
  \item \textbf{Solitary Kidney with Inadequate Function:} If the patient has only one kidney, or the contralateral kidney has insufficient function to sustain life, nephrectomy is absolutely contraindicated unless the patient is already on dialysis or dialysis is planned post-operatively.
  \item \textbf{Uncontrolled Coagulopathy:} Severe, uncorrectable bleeding disorders that significantly increase the risk of intraoperative or postoperative hemorrhage.
  \item \textbf{Severe Cardiopulmonary Disease:} Life-threatening cardiac or pulmonary conditions that make general anesthesia and major surgery prohibitively risky.
  \item \textbf{Extensive Metastatic Disease with Poor Performance Status:} For oncologic indications, if the cancer is widely metastatic and the patient's overall health (performance status) is very poor, surgery may offer no survival benefit and only increase morbidity.
\end{itemize}

\textbf{Relative Contraindications and Precautions:}
\begin{itemize}
  \item \textbf{Significant Comorbidities:} Conditions such as severe obesity, uncontrolled diabetes, or advanced liver disease increase surgical complexity and risk of complications. Careful optimization of these conditions is required.
  \item \textbf{Prior Abdominal Surgery:} Extensive adhesions from previous operations can make dissection more challenging and increase the risk of bowel injury, especially with minimally invasive approaches.
  \item \textbf{Large Tumor Size or Complex Anatomy:} Very large tumors, those with renal vein or inferior vena cava (IVC) thrombus, or those involving adjacent organs may necessitate an open approach or require specialized surgical expertise.
  \item \textbf{Active Systemic Infection:} Surgery should generally be postponed until any active systemic infection is treated and resolved.
  \item \textbf{Patient Refusal:} A patient's informed decision to decline surgery, after understanding all risks and benefits, is a contraindication.
\end{itemize}

\section*{Surgical Technique \& Procedure}
Nephrectomy is performed under general anesthesia, ensuring the patient is completely unconscious and pain-free throughout the procedure. Patient positioning varies based on the surgical approach: a lateral decubitus (flank) position is common for open or laparoscopic approaches, while a supine or modified flank position may be used for robotic surgery. The choice of approach (open, laparoscopic, or robotic) depends on tumor characteristics, patient comorbidities, and surgeon expertise.

A common approach, such as a laparoscopic radical nephrectomy, typically involves the following steps:
\begin{enumerate}
  \item \textbf{Anesthesia and Positioning:} General anesthesia is induced. The patient is typically placed in a modified lateral decubitus position, allowing access to the flank. The operating table is often flexed to widen the space between the iliac crest and costal margin.
  \item \textbf{Incision and Port Placement:} Several small incisions (typically 3-5, 5-12 mm in length) are made in the abdomen or flank. A Veress needle or an optical trocar is used to establish pneumoperitoneum by insufflating carbon dioxide gas, creating a working space. Trocars are then inserted through the incisions to accommodate the camera and surgical instruments.
  \item \textbf{Colon Mobilization:} For a left nephrectomy, the descending colon is mobilized medially by incising the line of Toldt. For a right nephrectomy, the ascending colon is mobilized. This exposes the retroperitoneal space and Gerota's fascia.
  \item \textbf{Kidney Dissection:} The kidney is identified within Gerota's fascia. The surrounding fat and attachments are carefully dissected using electrocautery or ultrasonic shears, freeing the kidney from its posterior and superior attachments. The adrenal gland is typically left intact during partial nephrectomy but removed with the kidney during radical nephrectomy for malignancy.
  \item \textbf{Hilar Dissection and Ligation:} The renal hilum, containing the renal artery and vein, is meticulously dissected. The renal artery is identified, isolated, and typically ligated first using vascular clips or an endoscopic stapler. The renal vein is then similarly ligated and divided. This sequence minimizes blood loss and prevents tumor dissemination.
  \item \textbf{Ureteral Division:} The ureter is identified, dissected, clipped or ligated, and divided close to the bladder.
  \item \textbf{Specimen Retrieval:} For radical nephrectomy, the entire kidney is freed and placed into an endoscopic retrieval bag. The bag is then extracted through one of the port sites, often slightly enlarged (e.g., a Pfannenstiel incision in the lower abdomen) to accommodate the specimen. For partial nephrectomy, after tumor excision and renorrhaphy, the kidney is inspected for hemostasis.
  \item \textbf{Closure:} The surgical field is inspected for bleeding. Drains may be placed if significant fluid accumulation is anticipated. The carbon dioxide gas is deflated, and the port sites are closed in layers, typically with absorbable sutures for the fascia and skin staples or sutures for the skin.
\end{enumerate}

\section*{Complications \& Risks}
As with any major surgical procedure, nephrectomy carries inherent risks, which can be broadly categorized into general surgical risks and procedure-specific complications.

\textbf{General Surgical Risks:}
\begin{itemize}
  \item \textbf{Anesthesia Complications:} These can include adverse reactions to anesthetic agents, respiratory depression, cardiac events (e.g., myocardial infarction, arrhythmia), stroke, or aspiration pneumonia.
  \item \textbf{Bleeding:} Intraoperative hemorrhage requiring blood transfusion, or postoperative hematoma formation.
  \item \textbf{Infection:} Surgical site infection (wound infection), urinary tract infection, pneumonia, or intra-abdominal abscess.
  \item \textbf{Deep Vein Thrombosis (DVT) and Pulmonary Embolism (PE):} Formation of blood clots in the legs that can travel to the lungs, a potentially life-threatening complication.
\end{itemize}

\textbf{Procedure-Specific Risks:}
\begin{itemize}
  \item \textbf{Acute Kidney Injury (AKI):} Especially a concern after partial nephrectomy due to ischemia-reperfusion injury, or in patients with pre-existing chronic kidney disease or a compromised contralateral kidney.
  \item \textbf{Damage to Adjacent Organs:}
    \begin{itemize}
      \item \textbf{Spleen:} Particularly during left nephrectomy, leading to hemorrhage or requiring splenectomy.
      \item \textbf{Liver:} During right nephrectomy, potentially causing bleeding or bile leak.
      \item \textbf{Bowel:} Injury to the colon or small intestine, leading to perforation, peritonitis, or fistula formation.
      \item \textbf{Pancreas:} Injury during left nephrectomy, potentially leading to pancreatitis or pancreatic fistula.
      \item \textbf{Diaphragm:} Injury can lead to pneumothorax (collapsed lung) or pleural effusion.
    \end{itemize}
  \item \textbf{Hernia Formation:} Incisional hernia at the site of an open incision or port-site hernia after laparoscopic/robotic surgery.
  \item \textbf{Urinary Leak:} Specific to partial nephrectomy, leakage of urine from the reconstructed collecting system, potentially forming a urinoma or fistula.
  \item \textbf{Chylous Ascites:} A rare complication resulting from injury to lymphatic vessels, leading to accumulation of lymphatic fluid in the abdominal cavity.
  \item \textbf{Nerve Injury:} Damage to peripheral nerves (e.g., ilioinguinal, genitofemoral) during dissection or positioning, leading to numbness or pain in the flank or groin.
  \item \textbf{Adrenal Insufficiency:} If the adrenal gland is removed (often with radical nephrectomy) and the contralateral adrenal gland is dysfunctional, or if both adrenal glands are removed.
  \item \textbf{Post-nephrectomy Syndrome:} Persistent flank pain or discomfort after kidney removal.
\end{itemize}

\section*{Postoperative Care \& Recovery}
Immediately after nephrectomy, the patient is transferred to the Post-Anesthesia Care Unit (PACU) for close monitoring. Vital signs, urine output, pain levels, and surgical site for bleeding are continuously assessed. Pain management is initiated, often with intravenous analgesics, patient-controlled analgesia (PCA), or epidural catheters. Early ambulation is encouraged to prevent complications like deep vein thrombosis and pneumonia.

Over the first 24-48 hours, the patient's condition is stabilized. Intravenous fluids are administered, and diet is gradually advanced from clear liquids to solid foods as tolerated. Drains, if placed, are monitored for output, and the Foley catheter is typically removed once the patient is mobile and able to void spontaneously. Nurses assist with deep breathing exercises and incentive spirometry to prevent pulmonary complications. The length of hospital stay varies, typically 2-4 days for minimally invasive approaches and 5-7 days for open surgery, depending on recovery progress and absence of complications. Over the first 1-2 weeks, patients are advised to restrict heavy lifting and strenuous activities, focusing on light ambulation and wound care. Full recovery, including return to normal activities and resolution of fatigue, can take several weeks to months, with gradual return to normal activities.

During the nephrectomy, the surgeon documents key intraoperative findings, including the size and location of the renal mass, its relationship to surrounding structures, any evidence of local invasion, and the appearance of the contralateral kidney. After removal, the resected kidney specimen is sent to pathology for definitive analysis. The pathology report is crucial, especially for oncologic cases, as it provides definitive diagnosis and staging information that guides prognosis and further treatment decisions. Key elements of the pathology report include the specific tumor type (e.g., clear cell renal cell carcinoma, papillary RCC, chromophobe RCC) or benign pathology, tumor size and grade (e.g., Fuhrman or ISUP grade), and crucially, the surgical margin status (negative or positive). It also details the presence or absence of lymphovascular invasion, involvement of adjacent structures (e.g., renal capsule, perirenal fat, adrenal gland, renal vein), and lymph node status if lymphadenectomy was performed. Based on these factors, a pathological TNM (Tumor, Node, Metastasis) stage is assigned.

A normal, successful postoperative course following nephrectomy is characterized by a predictable progression of healing and recovery. Patients can expect initial incisional pain, which is managed effectively with analgesics and gradually subsides over days to weeks. Return of bowel function, indicated by passing flatus and bowel movements, typically occurs within 1-3 days. The surgical wound should heal cleanly without signs of infection, with sutures or staples removed at the first follow-up visit. Patients will gradually regain strength and mobility, with fatigue being common in the initial weeks. For those undergoing partial nephrectomy, the remaining kidney should demonstrate stable function, while for radical nephrectomy, the solitary kidney should adequately compensate, maintaining stable serum creatinine levels. The goal is a progressive return to baseline activities and quality of life, free from major complications, within the expected recovery period.

Post-nephrectomy follow-up is crucial for monitoring recovery, assessing long-term renal function, and for cancer surveillance in oncologic cases.

Typical follow-up includes:
\begin{itemize}
  \item \textbf{Initial Postoperative Visit:} Scheduled 2-4 weeks after surgery for wound check, removal of sutures/staples, review of the final pathology report, and discussion of any adjuvant therapy if indicated.
  \item \textbf{Laboratory Tests:} Regular monitoring of serum creatinine and GFR to assess the function of the remaining kidney. This may be done at 1-2 months post-op, then every 6-12 months, or as clinically indicated.
  \item \textbf{Imaging for Cancer Surveillance:} For patients treated for renal cancer, periodic imaging (e.g., CT scan of the abdomen and pelvis, chest X-ray or CT chest) is performed to detect local recurrence or distant metastases. The frequency and type of imaging depend on the tumor stage and risk stratification, typically ranging from every 3-6 months initially, then annually for several years.
  \item \textbf{Activity Resumption:} Gradual increase in physical activity, with restrictions on heavy lifting (usually for 6-8 weeks) to allow for proper wound healing and prevent hernia formation.
  \item \textbf{Rehabilitation/Physical Therapy:} May be recommended for some patients to regain strength and mobility, especially after open surgery.
  \item \textbf{Lifestyle Modifications:} Counseling on maintaining a healthy lifestyle, including diet, hydration, and blood pressure control, to protect the remaining kidney.
  \item \textbf{Warning Signs:} Patients are educated on warning signs that warrant immediate medical attention, such as fever, chills, severe or worsening pain, redness or purulent drainage from the wound, persistent nausea or vomiting, decreased urine output, or swelling in the legs.
\end{itemize}
Adherence to these follow-up recommendations is vital for optimal long-term outcomes.

\section*{Outcomes \& Prognosis}
Nephrectomy is a commonly performed urological procedure, with its epidemiology largely driven by the incidence of renal cell carcinoma (RCC) and, to a lesser extent, benign renal conditions and living kidney donation. Renal cell carcinoma accounts for approximately 2-3\% of all adult malignancies, with an estimated 79,000 new cases diagnosed annually in the United States. The incidence of RCC has been steadily increasing, partly due to the widespread use of cross-sectional imaging, leading to the incidental discovery of smaller, asymptomatic renal masses. RCC typically affects individuals in their 60s and 70s and is more common in men than women, with a male-to-female ratio of approximately 1.5:1. While RCC is the primary indication, nephrectomy is also performed for benign conditions such as severe hydronephrosis, chronic pyelonephritis, and large angiomyolipomas, though these represent a smaller proportion of cases. Living kidney donation procedures contribute significantly to the overall number of nephrectomies performed, with thousands of donor nephrectomies conducted globally each year. The mortality rate for elective nephrectomy is generally low, typically less than 1-2\% in experienced centers, but can be higher for emergency procedures or in patients with significant comorbidities. Morbidity rates, encompassing various complications, range from 10-20\% depending on the surgical approach (open vs. minimally invasive), tumor complexity, and patient factors.

The outcomes and prognosis following nephrectomy are highly dependent on the underlying indication, the stage of disease (for cancer), the surgical approach, and the patient's overall health. For localized renal cell carcinoma, nephrectomy offers excellent oncologic control. Five-year survival rates for patients with localized RCC (Stage I and II) treated with radical or partial nephrectomy range from 70\% to over 90\%, with partial nephrectomy demonstrating comparable oncologic outcomes to radical nephrectomy for small renal masses (T1a and T1b). A key advantage of partial nephrectomy is its superior preservation of long-term renal function, which translates to a reduced risk of chronic kidney disease and associated cardiovascular morbidity, as highlighted by numerous comparative studies.

For patients undergoing nephrectomy for benign, non-functioning kidneys, the prognosis is generally excellent, with resolution of symptoms and no impact on overall survival, provided the contralateral kidney is healthy. Living kidney donors also experience excellent long-term outcomes, with no significant increase in the risk of end-stage renal disease or mortality compared to the general population. Prognostic factors for RCC include tumor stage (the most significant factor), Fuhrman/ISUP grade, presence of lymphovascular invasion, sarcomatoid differentiation, and performance status. Recurrence rates vary widely with stage, with higher-stage tumors having a greater risk of local recurrence or distant metastasis. Regular surveillance is crucial for early detection of recurrence, which can impact long-term survival.

\section*{References}
\begin{enumerate}
  \item MedlinePlus. Nephrectomy. U.S. National Library of Medicine; 2024. Available from: \url{https://medlineplus.gov/ency/article/002959.htm}
  \item Wein AJ, Kavoussi LR, Partin AW, Peters CA, editors. Campbell-Walsh-Wein Urology. 12th ed. Philadelphia, PA: Elsevier; 2021.
  \item Schwartz SI, Brunicardi FC, Andersen DK, Billiar TR, Hunter JG, Matthews JB, Pollock RE, editors. Schwartz's Principles of Surgery. 11th ed. New York, NY: McGraw-Hill Education; 2019.
  \item American Urological Association (AUA) Guidelines. Management of Small Renal Masses. 2021. Available from: \url{https://www.auanet.org/guidelines-and-quality/guidelines/small-renal-mass-guideline}
\end{enumerate}

\clearpage
